% ============================================================================
% Section 2.8 — Personal Financial Literacy
% State-specific: Minnesota, Oklahoma, Texas
% ============================================================================
\section{Personal Financial Literacy}

\begin{stepGoal}
\begin{itemize}
    \item Compare financial institutions and payment methods
    \item Understand how credit, debit, and interest affect your money
\end{itemize}
\end{stepGoal}

\begin{stepVocabBox}
    \vocabItem{Debit Card}{Spends money you already have in your account.}
    \vocabItem{Credit Card}{Borrows money you must pay back, often with interest.}
    \vocabItem{Interest Rate}{The percent a lender charges or a bank pays you.}
\end{stepVocabBox}

\begin{stepsCard}[How to Compare Payment Options]
    \stepItem{Identify the \textbf{payment method}: cash, debit, or credit.}
    \stepItem{Calculate any \textbf{extra cost} (interest, fees) using $I = Prt$ or the given rate.}
    \stepItem{Compare \textbf{total costs} to decide which option saves the most money.}
\end{stepsCard}

% --- Worked Example 1 ---
\begin{stepExample}{You buy a $\$200$ game console. Compare paying with a debit card vs.\ a credit card at $18\%$ annual rate paid after $1$ year.}
    \stepShow{1} Debit: you pay cash from your account. Credit: you borrow $\$200$.

    \stepShow{2} Credit interest: $I = 200 \times 0.18 \times 1 = \$36$.

    \stepShow{3} Debit total: $\$200$. Credit total: $200 + 36 = \$236$.
    \stepResult{Credit costs $\$36$ more than debit.}
\end{stepExample}

% --- Worked Example 2 ---
\begin{stepExample}{Bank~A pays $2\%$ interest on savings. Bank~B pays $3.5\%$. You deposit $\$1{,}000$ for $1$ year. Which earns more?}
    \stepShow{1} Both are savings accounts. Compare interest earned.

    \stepShow{2} Bank~A: $1{,}000 \times 0.02 \times 1 = \$20$. Bank~B: $1{,}000 \times 0.035 \times 1 = \$35$.

    \stepShow{3} Difference: $35 - 20 = \$15$.
    \stepResult{Bank~B earns $\$15$ more per year.}
\end{stepExample}

\mascotStepTip{Pay credit card bills on time every month to avoid extra interest charges!}

% ============================================================================
% PRACTICE
% ============================================================================
\newpage
\begin{stepPractice}{Financial Literacy Practice}
    \resetProblems

    \stepPracticeHeader[step1Color]{Identify the Payment Method}
    \prob You charge $\$150$ on a credit card at $15\%$ annual interest. What is the interest after $1$ year? \answerBlank[2cm]
    \answer{$\$22.50$}

    \stepPracticeHeader[step2Color]{Calculate Extra Costs}
    \begin{multicols}{2}
    \prob A checking account charges $\$5$/month. A credit union charges $\$2$/month. How much do you save per year at the credit union? \answerBlank[2cm]
    \answer{$\$36$}
    \prob A credit card charges $20\%$ on $\$300$ for $1$ year. What is the total you owe? \answerBlank[2cm]
    \answer{$\$360$}
    \end{multicols}

    \stepPracticeHeader[step3Color]{Compare and Decide}
    \trueOrFalse{A debit card lets you spend more than you have in your account.}
    \answerTF{False}

    \wordProblem{You want to buy a $\$500$ laptop. Option~A: pay cash now. Option~B: credit card at $16\%$ annual interest, paid after $1$ year. How much more does Option~B cost?}{\$}
    \answer{$\$80$}
\end{stepPractice}
